%% minimalist-manual.tex -- User manual of the "minimalist" Beamer theme.
%% This manual is itself a presentation built with the theme.
%% Build:  latexmk -pdf minimalist-manual.tex   (run twice for frame totals)
\documentclass[11pt,aspectratio=169,t]{beamer}

\usepackage[english]{babel}
\usetheme[footer=total,titlerule,sectionpages]{minimalist}

\tcbuselibrary{listings}
% "Code | Result" boxes used throughout this manual
\tcbset{
  demo base/.style={
    listing options={language={[LaTeX]TeX},basicstyle=\ttfamily\scriptsize,
      breaklines=true,backgroundcolor={},frame=none,xleftmargin=0pt},
    sharp corners, boxrule=0.4pt, colframe=black!55, colback=black!3,
    colbacklower=white, colupper=black!85, halign lower=flush left,
    left=4pt, right=4pt, top=3pt, bottom=3pt,
    before skip=0.3em, after skip=0.3em},
  demo/.style={demo base, listing side text,
    lefthand width=0.55\linewidth, sidebyside gap=6pt},
  demo stacked/.style={demo base, listing and text},
}

\title{Minimalist Beamer Theme}
\subtitle{User manual and example gallery}
\author{Version 1.0.0}
\institute{pdfLaTeX \textbullet\ XeLaTeX \textbullet\ LuaLaTeX}
\date{October 2026}

\begin{document}

\begin{frame}[c]\titlepage\end{frame}

\begin{frame}{Contents}
  \tableofcontents
\end{frame}

%% =====================================================================
\section{Getting started}

\begin{frame}{What is this theme?}
  A \concept{minimalist, grayscale} Beamer theme: white background, sans-serif
  text, no navigation symbols, and a small set of reusable building blocks.

  \begin{twocol}
    \concept{Layout}
    \begin{itemize}
      \item Left-aligned title page and frame titles
      \item Optional section divider pages
      \item \code{onecol}, \code{twocol} and \code{threecol} environments
    \end{itemize}
    \nextcol
    \concept{Content}
    \begin{itemize}
      \item Figures: TikZ, \code{\textbackslash flow}, images
      \item Tables: \code{booktabs}, \code{L}/\code{C}/\code{R}
      \item Boxes: \code{callout}, \code{statement}, \dots
      \item Code listings and math
    \end{itemize}
  \end{twocol}

  \begin{notebox}
    Documentation and commands are in English; your slides can be in any
    \code{babel} language (e.g.\ Spanish: ñ, á, ¿, ¡).
  \end{notebox}
\end{frame}

\begin{frame}[fragile]{Quick start}
\begin{lstlisting}[language={[LaTeX]TeX}]
\documentclass[11pt,aspectratio=169]{beamer}
\usepackage[spanish,es-tabla]{babel}   % or english, etc.
\usetheme[footer=total]{minimalist}

\title{My talk}
\author{Me}
\begin{document}
\frame{\titlepage}
\frame{\frametitle{Hello}
  Some text and a \concept{key term}.}
\end{document}
\end{lstlisting}
  Put \code{beamerthememinimalist.sty} next to your \code{.tex} file (or in
  your local \code{texmf} tree) and compile with
  \code{latexmk -pdf}.
\end{frame}

\begin{frame}{Theme options}
  \small
  Pass options to the theme: \code{\textbackslash usetheme[\textit{options}]\{minimalist\}}.

  \vspace{0.5em}
  \begin{tabular}{@{}l l l L{5.6cm}@{}}
    \toprule
    \textbf{Option} & \textbf{Values} & \textbf{Default} & \textbf{Effect}\\
    \midrule
    \code{accent}       & any color expr.        & \code{black!75} & Color of bullets, rules, box frames and titles\\
    \code{footer}       & \code{number}, \code{total}, \code{none} & \code{number} & Frame counter style (bottom right)\\
    \code{font}         & \code{helvetica}, \code{default} & \code{helvetica} & Helvetica clone, or Beamer's default sans font\\
    \code{titlerule}    & boolean                & \code{false}    & Thin rule under each frame title\\
    \code{sectionpages} & boolean                & \code{false}    & Automatic divider frame at each \code{\textbackslash section}\\
    \bottomrule
  \end{tabular}

  \vspace{0.6em}
  Example: \code{\textbackslash usetheme[accent=blue!60!black,footer=total,titlerule]\{minimalist\}}
\end{frame}

\begin{frame}[fragile]{Languages and TeX engines}
  \begin{twocol}[split=0.5]
    \concept{Languages}
    \begin{itemize}
      \item The theme never loads \code{babel}: you choose the language.
      \item Captions follow \code{babel}: \emph{Figure}, \emph{Figura}, \emph{Tabla}.
      \item With Spanish, \code{es-tabla} gives ``Tabla'' instead of ``Cuadro''.
      \item Accents and \code{ñ ¿ ¡} work out of the box (UTF-8).
    \end{itemize}
    \nextcol
    \concept{Engines}
    \begin{itemize}
      \item \textbf{pdfLaTeX}: T1 fonts and the Helvetica clone.
      \item \textbf{XeLaTeX / LuaLaTeX}: \code{fontspec} with TeX Gyre Heros.
      \item Detection is automatic (\code{iftex}).
      \item Listings accept accented characters under all engines.
    \end{itemize}
  \end{twocol}
\begin{lstlisting}[language={[LaTeX]TeX}]
\usepackage[spanish,es-tabla]{babel}  % load BEFORE \usetheme
\usetheme{minimalist}
\end{lstlisting}
\end{frame}

%% =====================================================================
\section{Page layout}

\begin{frame}{Title page and sections}
  The \code{\textbackslash titlepage} shows title, subtitle, author,
  institute and date, left aligned with a short accent rule.

  \vspace{0.4em}
  \begin{itemize}
    \item \code{\textbackslash section\{\ldots\}} marks a section; with the option
          \code{sectionpages} a divider frame is inserted automatically.
    \item Without that option, insert one manually:
          \code{\textbackslash begin\{frame\}\textbackslash sectionpage\textbackslash end\{frame\}}.
    \item \code{\textbackslash tableofcontents} uses numbered sections.
    \item \code{\textbackslash framesubtitle\{\ldots\}} adds a subtitle under the frame title.
  \end{itemize}
\end{frame}

\begin{frame}[fragile]{One column}
  By default content uses the full text width, which is \concept{one column}.
  For a narrower, centered measure use \code{onecol}:

\begin{tcblisting}{demo}
\begin{onecol}[width=0.7]
  A narrower text block, centered
  on the slide. Good for quotes,
  definitions or short statements.
\end{onecol}
\end{tcblisting}
  \code{width} is a fraction of the text width (default \code{0.8}).
\end{frame}

\begin{frame}{Two columns: result}
  \framesubtitle{Use \code{twocol} on the slides that need it}
  \begin{twocol}[split=0.45]
    \concept{Left column}
    \begin{itemize}
      \item Plain text, lists, formulas\ldots
      \item Any content works inside a column.
    \end{itemize}
    \nextcol
    \begin{callout}[title=Right column]
      Boxes, figures and tables fit too.
    \end{callout}
  \end{twocol}
  Content after \code{\textbackslash end\{twocol\}} goes back to the full width.
\end{frame}

\begin{frame}[fragile]{Two or three columns: source and keys}
\begin{lstlisting}[language={[LaTeX]TeX}]
\begin{twocol}[split=0.45]
  Left column content
\nextcol
  Right column content
\end{twocol}
% three columns: \begin{threecol} A \nextcol B \nextcol C \end{threecol}
\end{lstlisting}
  \small
  \begin{tabular}{@{}l l L{6.6cm}@{}}
    \toprule
    \textbf{Key} & \textbf{Default} & \textbf{Meaning}\\
    \midrule
    \code{split} & \code{0.5}  & Left column share of the width (\code{twocol} only)\\
    \code{gap}   & \code{0.04} / \code{0.06}$^*$ & Fraction of the width left empty between columns\\
    \code{align} & \code{T}    & Vertical alignment: \code{T}, \code{c} or \code{b}\\
    \bottomrule
  \end{tabular}
  \par\smallskip{\footnotesize $^*$ default for \code{twocol} / \code{threecol}.}
\end{frame}

%% =====================================================================
\section{Building blocks}

\begin{frame}[fragile]{Text helpers}
\begin{tcblisting}{demo}
\concept{Key term} in bold,
\code{inline\_code()} in monospace,
\alert{alert} for emphasis.

\flowline{Design, Code, Test, Integrate}

Single arrow: A \arr\ B
\end{tcblisting}
  Standard \code{itemize}, \code{enumerate} and \code{description} keep
  Beamer's behavior, tinted with the accent color.
\end{frame}

\begin{frame}[fragile]{Figures: process flows with \texttt{\textbackslash flow}}
  \code{\textbackslash flow\{A, B, C\}} draws a centered chain of nodes joined by
  arrows and shrinks it to fit the line. Use \code{\textbackslash\textbackslash} for line breaks.

\begin{tcblisting}{demo stacked}
\flow{Requirements, Domain\\model, Design, Code}
\end{tcblisting}
  Pass TikZ options in brackets:
  \code{\textbackslash flow[node distance=1cm]\{\ldots\}}.
\end{frame}

\begin{frame}[fragile]{Figures: custom TikZ with theme styles}
  Styles: \code{mn node}, \code{mn arrow}, \code{mn dashed arrow}, \code{mn circle}.

\begin{tcblisting}{demo stacked, listing options={basicstyle=\ttfamily\tiny}}
\begin{tikzpicture}[mn circle/.append style={minimum size=1.5cm}]
  \node[mn circle] (r) {RED};
  \node[mn circle, right=1.4cm of r] (g) {GREEN};
  \node[mn circle, right=1.4cm of g] (f) {REFACTOR};
  \draw[mn arrow] (r) -- (g);  \draw[mn arrow] (g) -- (f);
  \draw[mn dashed arrow, bend left=30] (f) to (r);
\end{tikzpicture}
\end{tcblisting}
\end{frame}

\begin{frame}[fragile]{Figures: images and captions}
  \begin{twocol}[split=0.5,align=c]
\begin{lstlisting}[language={[LaTeX]TeX}]
\begin{figure}
  \centering
  \includegraphics[width=0.6\linewidth]%
    {example-image-a}
  \caption{An example image.}
\end{figure}
\end{lstlisting}
    Captions are numbered and follow the \code{babel} language.
    \nextcol
    \begin{figure}
      \centering
      \includegraphics[width=0.6\linewidth]{example-image-a}
      \caption{An example image.}
    \end{figure}
  \end{twocol}
\end{frame}

\begin{frame}[fragile]{Tables}
  Use \code{booktabs} rules and the column types \code{L\{w\}}, \code{C\{w\}}, \code{R\{w\}}
  (left, centered, right; fixed width, wrapped text).

\begin{tcblisting}{demo, lefthand width=0.58\linewidth}
\begin{table}
  \centering
  \begin{tabular}{@{}L{1.7cm}C{1cm}R{1.1cm}@{}}
    \toprule
    \textbf{Item} & \textbf{Qty} & \textbf{Price}\\
    \midrule
    Pencil & 3 & 1.20\\
    Notebook & 2 & 4.50\\
    \bottomrule
  \end{tabular}
  \caption{A table.}
\end{table}
\end{tcblisting}
\end{frame}

\begin{frame}[fragile]{Boxes: \texttt{callout} and \texttt{statement}}
\begin{tcblisting}{demo}
\begin{callout}
  Outlined box for a key idea.
\end{callout}

\begin{callout}[title=With a title]
  Any tcolorbox key can be passed.
\end{callout}

\begin{statement}
  A centered, emphasized sentence.
\end{statement}
\end{tcblisting}
\end{frame}

\begin{frame}[fragile]{Boxes: \texttt{notebox} and Beamer blocks}
\begin{tcblisting}{demo}
\begin{notebox}[title=Note]
  Soft gray note with a left bar.
\end{notebox}

\begin{block}{Block}
  Standard Beamer block.
\end{block}

\begin{alertblock}{Alert}
  Use sparingly.
\end{alertblock}
\end{tcblisting}
\end{frame}

\begin{frame}[fragile]{Code listings}
  \alert{Frames with listings need the \code{[fragile]} option.}
  The \code{minimalist} listings style is already active:

\begin{lstlisting}[language=Python]
def peso_total(envio):
    # Suma el peso de cada línea de suministro
    return sum(l.peso * l.cantidad for l in envio.lineas)
\end{lstlisting}
  Inline: \lstinline{assert peso_total(e) <= v.capacidad}.
  Choose another language with \code{[language=Java]}, or set it for the
  whole deck with \code{\textbackslash lstset\{language=Python\}}.
\end{frame}

\begin{frame}[fragile]{Math, definitions and theorems}
\begin{tcblisting}{demo}
\begin{definition}
  A set $S$ is \emph{finite} if\ldots
\end{definition}
\begin{theorem}
  $\displaystyle\sum_{i=1}^{n} i = \frac{n(n+1)}{2}$
\end{theorem}
\end{tcblisting}
  The names (``Definition'', ``Teorema'', \ldots) are translated by \code{babel}.
\end{frame}

%% =====================================================================
\section{Reference}

\begin{frame}{Command reference}
  \footnotesize
  \begin{tabular}{@{}L{5.6cm}L{7.8cm}@{}}
    \toprule
    \textbf{Command / environment} & \textbf{Purpose}\\
    \midrule
    \code{twocol[split,gap,align]}, \code{\textbackslash nextcol} & Two columns; \code{\textbackslash nextcol} starts the second one\\
    \code{threecol[gap]}, \code{\textbackslash nextcol} & Three equal columns (same \code{\textbackslash nextcol})\\
    \code{onecol[width]}      & Centered single column of a given fraction of the width\\
    \code{callout[opts]}      & Outlined box, optional \code{title=}\\
    \code{statement[opts]}    & Centered, emphasized box\\
    \code{notebox[opts]}      & Gray note with left bar\\
    \code{\textbackslash concept\{\}}, \code{\textbackslash code\{\}}, \code{\textbackslash arr} & Bold term, monospace, arrow\\
    \code{\textbackslash flowline\{A,B,C\}} & Inline chain  A \arr\ B \arr\ C\\
    \code{\textbackslash flow[opts]\{A,B,C\}} & TikZ node chain joined by arrows\\
    \code{mn node}, \code{mn arrow}, \code{mn dashed arrow}, \code{mn circle} & TikZ styles\\
    \code{L\{w\}}, \code{C\{w\}}, \code{R\{w\}} & Fixed-width table columns\\
    \code{lstlisting} (style \code{minimalist}) & Code listings (fragile frames)\\
    \bottomrule
  \end{tabular}
\end{frame}

\begin{frame}{Customization}
  \begin{itemize}\small
    \item \concept{Aspect ratio:} \code{\textbackslash documentclass[aspectratio=43]\{beamer\}}.
    \item \concept{Top alignment:} class option \code{t}; use \code{[c]} on a single frame to center it.
    \item \concept{Accent color:} \code{\textbackslash usetheme[accent=blue!50!black]\{minimalist\}}.
    \item \concept{Margins:} \code{\textbackslash setbeamersize\{text margin left=1.2cm\}}.
    \item \concept{Node look:} \code{\textbackslash tikzset\{mn node/.append style=\{fill=black!5\}\}}.
    \item \concept{Boxes:} \code{\textbackslash tcbset\{mn base/.append style=\{arc=3pt\}\}}.
    \item \concept{Table spacing:} \code{\textbackslash renewcommand\{\textbackslash arraystretch\}\{1.2\}}.
  \end{itemize}
  \vspace{0.4em}
  \begin{callout}
    Override these \emph{after} \code{\textbackslash usetheme}.
  \end{callout}
\end{frame}

\begin{frame}{Migrating from the original ad-hoc template}
  \small
  \begin{tabular}{@{}L{5.4cm}C{0.8cm}L{6.6cm}@{}}
    \toprule
    \textbf{Original} & & \textbf{Theme}\\
    \midrule
    \code{\textbackslash concepto\{\ldots\}} & \arr & \code{\textbackslash concept\{\ldots\}}\\
    \code{\textbackslash caja\{\textbackslash centering \ldots\}} & \arr & \code{statement} or \code{callout}\\
    \code{columns} + \code{\textbackslash column\{0.47\textbackslash linewidth\}} & \arr & \code{twocol[split=0.5]} + \code{\textbackslash nextcol}\\
    TikZ chain with \code{every node/.style} & \arr & \code{\textbackslash flow\{A, B, C\}} or \code{mn node}\\
    \code{p\{4cm\}} with \code{\textbackslash raggedright} & \arr & \code{L\{4cm\}}\\
    Preamble with fonts, colors, footline & \arr & \code{\textbackslash usetheme\{minimalist\}}\\
    \bottomrule
  \end{tabular}
\end{frame}

\end{document}
